\section{The Released Framework: Verification and Reproducibility}
\label{sec:artifact}

The pipeline is released as a dependency-free Python package (standard library
only; Python $\geq 3.12$). Its purpose in this paper is not to report results
but to make the protocol \emph{executable and checkable}: every stage below is
runnable, and every structural claim in Appendix~\ref{app:claims} names the test
that establishes it.

\subsection{Four verification gates}
\label{sec:gates}

Every emitted task directory must pass all four gates before it can be
packaged.

\begin{enumerate}
  \item \textbf{Strict schema.} An embedded JSON Schema with
        \texttt{additionalProperties: false} at every level. This is what
        catches \harbor's silent key-drop (claim~\cid{C1}).
  \item \textbf{Layout sanity.} The files \texttt{harbor run} actually needs,
        in the places it expects them.
  \item \textbf{Upstream cross-check.} Optional: the task directory is loaded
        with the real \harbor\ models when \harbor\ is installed. When it is
        absent the gate reports \textsc{Unverified}---a third state, distinct
        from pass and fail, so an environment can never be credited as verified
        on a machine that could not check it (claim~\cid{C2}).
  \item \textbf{Hygiene.} LF line endings and no path escaping the task
        directory. A stray CRLF turns a \texttt{tests/test.sh} into a syntax
        error inside a Linux container, and nothing upstream reports that
        clearly.
\end{enumerate}

\subsection{Resumability and provenance}

Stage progress is persisted per paper after every stage, so a crashed run
resumes and a resumed run issues \emph{zero} model calls (claim~\cid{C4}). Each
packaged environment carries its source paper's identifier and licence in
\texttt{task.toml} provenance, and the set-level manifest records the metric and
verification status of every environment.

\subsection{Offline reproducibility}

The full pipeline and all four protocol stages have offline paths: a replay
backend (canned responses), a dry-run backend (schema-valid synthetic output),
and manifest-driven scoring for Stage~1. Consequently the entire surface is
exercised in continuous integration without a model, a container runtime, or
network access (claim~\cid{C8}), and a reviewer can reproduce any run
deterministically.

\subsection{Limitations of the artifact}
\label{app:limits}

\begin{itemize}
  \item No environment has been constructed from a live paper and no model has
        been run in the released version; all empirical claims are
        \texttt{planned}.
  \item Variance decomposition is reliable only on complete model $\times$
        harness grids (\S\ref{sec:stage3}).
  \item The n-gram contamination probe can miss a verbatim leak when a probe
        text is shorter than the n-gram window, because a short probe's single
        gram cannot match an 8-gram of the training text. Reported overlap is
        therefore a lower bound.
  \item Harness deltas treat an absent key as $0.0$, so ``did not participate''
        can be read as ``regressed''.
  \item Stage~1 summaries do not yet expose the review and probe denominators,
        so the reviewed-only denominator of the expert approval rate is not
        auditable from the report alone.
\end{itemize}

\noindent These are tracked in the repository and are re-checked before any
number is transcribed into this manuscript.
